
@@M0@@, on a
@@M1@@
où @@M2@@ et @@M3@@
sont les morphismes structuraux de @@M4@@ et @@M5@@. Comme @@M6@@, on a:
@@M7@@
d'après IV 2.12.4.

D'après (3.4.1), on a donc :
@@M8@@
Soient @@M9@@ et @@M10@@ les classes des faisceaux inversibles canoniques sur @@M11@@ et @@M12@@. Alors, d'après VI 1.1, @@M13@@ (resp. @@M14@@) admet pour base sur @@M15@@ les éléments @@M16@@ pour @@M17@@ (resp. @@M18@@ pour @@M19@@). La formule de projection donne :
@@M20@@
Si @@M21@@ est l'idéal de @@M22@@ dans @@M23@@, @@M24@@ est un idéal inversible, et on a @@M25@@. Par ailleurs, @@M26@@ est défini par l'idéal @@M27@@ de @@M28@@ noyau de @@M29@@, qui est engendré par un sous-Module de @@M30@@ isomorphe à @@M31@@; par suite, @@M32@@. L'image de @@M33@@ est donc le sous-@@M34@@-module engendré par les @@M35@@ pour @@M36@@. L'homomorphisme @@M37@@ est donc un isomorphisme, et par suite, il en est de même pour l'homomorphisme @@M38@@.

\subsection*{4. Fibrés principaux sous les groupes @@M39@@}

\textbf{Proposition 4.1.} \uline{Soient} @@M40@@ \uline{un schéma noethérien}, @@M41@@ \uline{le} @@M42@@-\uline{groupe} @@M43@@, @@M44@@ \uline{un torseur sous} @@M45@@, \uline{et} @@M46@@ \uline{le} @@M47@@\uline{-Module localement libre de rang} @@M48@@ \uline{correspondant à} @@M49@@. \uline{On note} @@M50@@ \uline{la classe de} @@M51@@ \uline{dans} @@M52@@, \uline{et} @@M53@@ \uline{le} @@M54@@-\uline{idéal engendré}
